% Glossary section for the chartbook
% Include in your main document with % Glossary section for the chartbook
% Include in your main document with % Glossary section for the chartbook
% Include in your main document with % Glossary section for the chartbook
% Include in your main document with \input{glossary}
% Requires: \usepackage{multicol}

\section*{Glossary}
\addcontentsline{toc}{section}{Glossary}

\begin{multicols}{2}
\small
\setlength{\parindent}{0pt}
\setlength{\parskip}{0.5\baselineskip}

\textbf{aggregate demand} --- The total demand for all goods and services in an economy at a given time. It combines spending by consumers, businesses, the government, and foreign buyers.

\textbf{annualized / annualized rate} --- A way of expressing a short-term change (such as monthly or quarterly) as if it continued for a full year. For example, if GDP grew 0.5\% in one quarter, the annualized rate would be about 2\%---the growth you'd see if that pace held for four quarters.

\textbf{average hourly earnings (AHE)} --- The average amount workers earn per hour, calculated from employer payroll data. This measure helps track wage growth but can be misleading when the mix of jobs changes (e.g., if low-wage workers lose jobs, the average rises even without anyone getting a raise).

\textbf{BEA (Bureau of Economic Analysis)} --- A federal agency within the Commerce Department that produces official statistics on the U.S. economy, including GDP, personal income, and international trade data.

\textbf{business cycle} --- The recurring pattern of expansion and contraction in economic activity over time, typically measured by changes in real GDP. A full cycle includes periods of growth followed by recession.

\textbf{capital flows} --- The movement of money across national borders for investment purposes, such as foreigners buying U.S. stocks or Americans investing in overseas factories.

\textbf{capital gains} --- The profit earned when an asset (like a stock or house) is sold for more than its purchase price. Capital gains are typically taxed at lower rates than regular income.

\textbf{composition effect} --- A statistical phenomenon where changes in average wages can be influenced by shifts in who is working, not just changes in individual pay. For example, if many low-wage workers lose their jobs during a recession, the average wage may rise even though no one received a raise---simply because the remaining workforce skews toward higher earners.

\textbf{consumer spending} --- The total amount of money households spend on goods and services. It is the largest component of GDP in the United States, typically accounting for about two-thirds of economic activity.

\textbf{contribution to GDP growth} --- The portion of overall GDP growth attributable to a specific component (like consumer spending or business investment). If GDP grew 3\% and consumer spending contributed 2 percentage points, that means consumer spending alone would have produced 2\% growth.

\textbf{core inflation} --- A measure of inflation that excludes food and energy prices, which tend to be volatile. Economists use core inflation to identify underlying price trends without the noise of short-term swings in gas or grocery prices.

\textbf{corporate profits} --- The earnings companies have left after paying their costs, including wages, materials, and interest. This measure helps gauge business health and can signal future hiring or investment.

\textbf{CPI (Consumer Price Index)} --- A measure of the average change over time in the prices paid by consumers for a basket of common goods and services, such as food, housing, and transportation. It is one of the most widely used measures of inflation.

\textbf{CPS (Current Population Survey)} --- A monthly survey of about 60,000 households conducted by the Census Bureau and Bureau of Labor Statistics. It provides the official unemployment rate and other labor market data based on asking people directly about their work status.

\textbf{current account (international)} --- A measure of a country's transactions with the rest of the world, including trade in goods and services, investment income, and transfers like remittances. When outflows exceed inflows, the country runs a current account deficit. The U.S. current account deficit is driven primarily by its trade deficit, but also reflects income paid to foreign investors.

\textbf{debt-to-income ratio} --- The percentage of a person's or household's income that goes toward paying debts. A ratio of 30\% means 30 cents of every dollar earned goes to debt payments.

\textbf{deflated} --- Adjusted to remove the effects of inflation, allowing for meaningful comparisons of economic values across different time periods. Also referred to as being expressed in ``real'' or ``constant'' dollars.

\textbf{depreciation} --- The decline in value of an asset over time due to wear, age, or obsolescence. In the context of currency, it refers to a decrease in a currency's value relative to other currencies.

\textbf{discretionary spending} --- Government spending that Congress decides on each year through the appropriations process, as opposed to mandatory spending (like Social Security) that continues automatically. It includes defense, education, and transportation.

\textbf{disposable personal income} --- The amount of money households have available for spending or saving after paying taxes. It is calculated as total personal income minus personal taxes.

\textbf{durable goods} --- Products designed to last three years or more, such as cars, appliances, and furniture. Spending on durable goods tends to fluctuate more than spending on other goods during the business cycle.

\textbf{ECI (Employment Cost Index)} --- A quarterly measure of how much employers pay for labor, including both wages and benefits. Unlike average hourly earnings, the ECI holds the mix of jobs constant, giving a cleaner picture of whether compensation is actually rising.

\textbf{economic expansion} --- A period when the economy is growing---GDP is rising, jobs are being added, and business activity is increasing. Expansions end when a recession begins.

\textbf{economic growth} --- An increase in the production of goods and services in an economy over time, typically measured as the percentage change in real GDP.

\textbf{employment rate} --- The percentage of the working-age population that is employed. It is calculated by dividing the number of employed people by the total working-age population.

\textbf{employment-to-population ratio} --- The percentage of the working-age population (usually ages 16 and older) that has a job. Unlike the unemployment rate, this measure captures people who have stopped looking for work.

\textbf{equity} --- Ownership value in an asset after subtracting what is owed. For a home worth \$400,000 with a \$250,000 mortgage, the equity is \$150,000. For stocks, equity represents ownership shares in a company.

\textbf{establishment survey} --- A monthly survey of about 670,000 business locations that provides data on jobs, hours, and earnings. It produces the widely reported nonfarm payrolls number and is considered more reliable for job counts than the household survey.

\textbf{exchange rate} --- The price of one country's currency expressed in terms of another country's currency. For example, if one U.S. dollar equals 0.85 euros, that is the dollar-to-euro exchange rate.

\textbf{Fed Funds Rate} --- The interest rate banks charge each other for overnight loans. The Federal Reserve uses this rate as its main tool for influencing the economy---lowering it to encourage borrowing and spending, raising it to cool inflation.

\textbf{Federal Reserve} --- The central bank of the United States, responsible for conducting monetary policy, supervising banks, and maintaining financial stability. Often called ``the Fed.''

\textbf{financial accounts} --- A comprehensive record of the assets, liabilities, and net worth of all sectors in the economy---households, businesses, government, and the rest of the world. These accounts track who owns what and who owes what to whom. Also called the ``flow of funds.''

\textbf{fiscal stimulus} --- Government actions to boost the economy through increased spending or tax cuts. During recessions, stimulus puts money in people's pockets to encourage spending and support businesses.

\textbf{fiscal year} --- A 12-month period used for government accounting and budgeting purposes. The U.S. federal fiscal year runs from October 1 through September 30.

\textbf{FOMC (Federal Open Market Committee)} --- The Federal Reserve committee that sets interest rate policy. It meets eight times a year and includes the Fed's Board of Governors plus regional Fed bank presidents.

\textbf{GDP (Gross Domestic Product)} --- The total market value of all goods and services produced within a country during a specific period, usually a year or quarter. It is the most commonly used measure of an economy's size.

\textbf{GDPNow} --- A real-time estimate of U.S. economic growth published by the Federal Reserve Bank of Atlanta. Unlike official GDP figures, which are released quarterly with a delay, GDPNow updates frequently as new economic data becomes available, giving an early read on how the economy is performing.

\textbf{government deficit} --- The amount by which government spending exceeds government revenue in a given period. When the government collects more than it spends, it has a surplus.

\textbf{Great Recession} --- The severe economic downturn from December 2007 to June 2009, triggered by the housing market collapse and financial crisis. It was the worst U.S. recession since the 1930s, with unemployment peaking at 10\%.

\textbf{gross domestic income (GDI)} --- The total income earned producing goods and services in a country---wages, profits, and other income combined. In theory, GDI should equal GDP since every dollar spent becomes someone's income, but measurement differences cause small gaps.

\textbf{home equity} --- The portion of a home's value that the owner actually owns---the market value minus the remaining mortgage balance. Home equity often represents the largest asset for middle-class families.

\textbf{homeownership rate} --- The percentage of households that own their home rather than rent. The U.S. rate has historically hovered around 65\%.

\textbf{household survey} --- The monthly survey of individuals and families (the Current Population Survey) that produces the official unemployment rate by asking people directly whether they are working or looking for work.

\textbf{income inequality} --- The extent to which income is distributed unevenly among a population. Higher inequality means a larger gap between what high earners and low earners receive.

\textbf{inflation} --- A general increase in prices across the economy over time, which reduces the purchasing power of money. When inflation occurs, each dollar buys fewer goods and services than before.

\textbf{inflation expectations} --- What people and businesses think inflation will be in the future. These expectations matter because they influence wage negotiations, pricing decisions, and interest rates.

\textbf{inflation rate} --- The percentage change in the price level over a specific period, usually measured annually. It indicates how quickly prices are rising or falling.

\textbf{initial claims} --- The number of people filing for unemployment insurance benefits for the first time each week. A leading indicator of labor market health---rising claims often signal upcoming job losses.

\textbf{insured unemployed} --- People who are currently receiving unemployment insurance benefits. This count is smaller than total unemployment because not everyone who loses a job qualifies for benefits, and benefits eventually expire for those who remain unemployed.

\textbf{inverted yield curve} --- When short-term interest rates are higher than long-term rates, the opposite of normal. An inverted yield curve has preceded every U.S. recession in recent decades, making it a closely watched warning sign.

\textbf{job openings} --- Positions that employers are actively trying to fill. High job openings relative to unemployed workers indicate a tight labor market where workers have bargaining power.

\textbf{JOLTS (Job Openings and Labor Turnover Survey)} --- A monthly survey tracking job openings, hiring, and separations (quits, layoffs, and other departures). It provides insight into labor market dynamics beyond simple job counts.

\textbf{labor force} --- The total number of people who are either employed or actively seeking employment. It does not include people who are not working and not looking for work.

\textbf{labor force participation rate} --- The percentage of the working-age population that is in the labor force (either employed or actively seeking work). It indicates how many people are engaged in or available for work.

\textbf{labor productivity} --- The amount of goods and services produced per hour of work. Higher productivity means workers are producing more output in the same amount of time.

\textbf{labor share of income} --- The portion of national income that goes to workers as wages and benefits, as opposed to capital owners as profits. A declining labor share means workers are capturing less of the economy's gains.

\textbf{liabilities} --- What a person, company, or government owes---debts and other financial obligations. On a balance sheet, assets minus liabilities equals net worth.

\textbf{liquidity} --- How easily an asset can be converted to cash without losing value. A savings account is highly liquid; a house is not, since selling takes time and involves costs.

\textbf{M2} --- A measure of the money supply that includes cash, checking accounts, savings accounts, and other easily accessible funds. The Federal Reserve tracks M2 to gauge how much money is available for spending in the economy.

\textbf{market capitalization} --- The total value of a company's outstanding stock, calculated by multiplying share price by number of shares. A company with 1 million shares at \$100 each has a market cap of \$100 million.

\textbf{means-tested benefits} --- Government programs that only provide assistance to people below certain income or asset thresholds, such as Medicaid, food stamps (SNAP), and housing assistance.

\textbf{median} --- The middle value in a set of numbers arranged in order. Half of all values fall above the median and half fall below. Unlike an average, the median is not skewed by extremely high or low values.

\textbf{monetary policy} --- Actions taken by the Federal Reserve to influence the availability and cost of money and credit in the economy. The Fed's primary tools include setting interest rates and buying or selling securities.

\textbf{net worth} --- The total value of assets (such as homes, savings, and investments) minus total liabilities (such as mortgages and debts). It represents overall financial wealth.

\textbf{NIPA (National Income and Product Accounts)} --- The official accounting system for the U.S. economy, maintained by the BEA. It includes GDP, national income, personal income, and related measures.

\textbf{nominal GDP} --- The total value of goods and services produced in an economy measured at current prices, without adjusting for inflation. Comparing nominal GDP across years can be misleading because price changes are included.

\textbf{nonfarm payrolls} --- The total number of paid workers in the U.S. economy, excluding farm workers, private household employees, and nonprofit workers. The monthly change in nonfarm payrolls is the most-watched indicator of job growth.

\textbf{owners' equivalent rent} --- A component of the inflation measure that estimates what homeowners would pay to rent their own homes. It's used because the costs of owning (like maintenance) are considered consumption, while the home purchase itself is investment.

\textbf{payroll taxes} --- Taxes deducted from workers' paychecks to fund Social Security and Medicare. Employers also pay a matching amount, making the combined rate about 15\% of wages up to certain limits.

\textbf{PCE (Personal Consumption Expenditures)} --- A measure of consumer spending on goods and services that is used to track price changes and calculate inflation. The Federal Reserve prefers the PCE price index over the CPI for setting monetary policy.

\textbf{percentage point} --- A unit for describing the difference between two percentages. If unemployment falls from 5\% to 4\%, it dropped by 1 percentage point (not 1\%, which would be a much smaller change to 4.95\%).

\textbf{percentile} --- A value indicating the percentage of a distribution that falls below it. For example, a household at the 75th percentile of income earns more than 75\% of all households.

\textbf{personal saving rate} --- The percentage of after-tax income that households save rather than spend. A rate of 5\% means households save 5 cents of every dollar of disposable income.

\textbf{poverty threshold} --- The income level below which a family is officially considered poor. The thresholds vary by family size and are updated yearly for inflation---in 2024, about \$31,000 for a family of four.

\textbf{PPI (Producer Price Index)} --- A measure of inflation at the wholesale level, tracking prices received by producers for their output. It's often seen as a leading indicator of consumer inflation.

\textbf{producer prices} --- The prices businesses receive for their goods and services at various stages of production. Changes in producer prices can signal future changes in consumer prices.

\textbf{productivity-pay gap} --- The growing divide between how much workers produce per hour and how much they are paid. Since the 1970s, productivity in the U.S. has risen significantly faster than typical worker compensation, meaning the economic gains from increased efficiency have not been evenly shared with workers.

\textbf{public debt} --- The total amount of money the government owes to creditors, accumulated from past deficits. Also called the national debt or government debt.

\textbf{quantitative easing} --- A Federal Reserve policy of buying large amounts of bonds to push down long-term interest rates and stimulate the economy when short-term rates are already near zero. The Fed used this tool extensively after the 2008 financial crisis and during the COVID-19 pandemic.

\textbf{quintile / income quintile} --- One of five equal groups (each representing 20\%) when a population is ranked from lowest to highest. The bottom quintile is the poorest 20\%; the top quintile is the richest 20\%. Often used to describe income distribution.

\textbf{real GDP} --- The value of goods and services produced in an economy, adjusted for inflation to reflect changes in actual output rather than price changes. It allows for meaningful comparisons across time.

\textbf{real interest rates} --- Interest rates adjusted for inflation, representing the true cost of borrowing or return on saving. If a bond pays 5\% and inflation is 3\%, the real interest rate is about 2\%.

\textbf{recession} --- A significant decline in economic activity lasting more than a few months, typically visible in falling real GDP, rising unemployment, and reduced consumer spending. In the United States, the National Bureau of Economic Research (NBER) officially determines when recessions begin and end.

\textbf{residential investment} --- Spending on housing, including new home construction, home improvements, and brokers' fees. It's a component of GDP that tends to lead the broader economy---declining before recessions and recovering before expansions.

\textbf{S\&P 500} --- A stock market index tracking 500 large U.S. companies, widely used as a benchmark for overall stock market performance. It represents about 80\% of U.S. stock market value.

\textbf{Sahm Rule} --- A recession indicator that signals an economic downturn when the three-month average unemployment rate rises at least 0.5 percentage points above its lowest point in the previous 12 months. Named after economist Claudia Sahm, this rule has reliably identified every U.S. recession since 1970.

\textbf{seasonally adjusted} --- Data that has been modified to remove predictable seasonal patterns, such as holiday shopping or summer hiring. This adjustment makes it easier to identify underlying economic trends.

\textbf{sectoral balances} --- The financial positions of the major parts of the economy---households, businesses, government, and the foreign sector---showing whether each is spending more than it earns (deficit) or earning more than it spends (surplus). Because one sector's deficit must be matched by surpluses elsewhere, these balances always sum to zero across the entire economy.

\textbf{shelter (as CPI category)} --- The housing component of the Consumer Price Index, including rent and owners' equivalent rent. Shelter is the largest single category in CPI, making up about one-third of the index.

\textbf{trade balance / trade deficit / net exports} --- The difference between what a country exports and imports. When exports exceed imports, the country has a trade surplus; when imports exceed exports, a trade deficit. In GDP calculations, this same measure is called ``net exports.'' The U.S. has run a trade deficit since the 1970s, meaning Americans buy more from abroad than they sell.

\textbf{transfer payments} --- Government payments to individuals that are not in exchange for goods or services, such as Social Security, unemployment benefits, and food assistance. Transfers redistribute income but don't directly produce GDP.

\textbf{Treasury yields} --- The interest rates on U.S. government bonds. Yields vary by maturity---short-term bills might yield 4\% while 30-year bonds yield 5\%---and serve as benchmarks for other interest rates throughout the economy.

\textbf{U3 unemployment rate} --- The official unemployment rate, measuring the percentage of the labor force that is jobless and actively looking for work. It does not count people who have given up searching or those working part-time who want full-time jobs.

\textbf{U6 unemployment rate} --- A broader measure of unemployment that includes the officially unemployed plus discouraged workers (who stopped looking) and part-time workers who want full-time work. U6 is typically several percentage points higher than the official rate.

\textbf{unemployment rate} --- The percentage of the labor force that is jobless and actively seeking work. It is one of the most closely watched indicators of economic health, reported monthly.

\textbf{unit labor costs} --- The labor cost to produce one unit of output, calculated as total labor compensation divided by total output. Rising unit labor costs can signal inflation pressure; falling costs suggest productivity gains.

\textbf{value added} --- The increase in value a business creates over and above the cost of inputs it purchases. A furniture maker's value added is the difference between what it pays for wood and what it sells furniture for. GDP is essentially the sum of value added across all businesses.

\textbf{yield} --- The income earned from an investment, usually expressed as an annual percentage of the investment's value. For bonds, the yield represents the return an investor receives from interest payments.

\textbf{yield curve} --- A chart showing interest rates on bonds of different maturities, typically Treasury securities. Normally, longer-term bonds pay higher rates; when short-term rates exceed long-term rates, the curve is inverted---a recession warning signal.

\end{multicols}

% Requires: \usepackage{multicol}

\section*{Glossary}
\addcontentsline{toc}{section}{Glossary}

\begin{multicols}{2}
\small
\setlength{\parindent}{0pt}
\setlength{\parskip}{0.5\baselineskip}

\textbf{aggregate demand} --- The total demand for all goods and services in an economy at a given time. It combines spending by consumers, businesses, the government, and foreign buyers.

\textbf{annualized / annualized rate} --- A way of expressing a short-term change (such as monthly or quarterly) as if it continued for a full year. For example, if GDP grew 0.5\% in one quarter, the annualized rate would be about 2\%---the growth you'd see if that pace held for four quarters.

\textbf{average hourly earnings (AHE)} --- The average amount workers earn per hour, calculated from employer payroll data. This measure helps track wage growth but can be misleading when the mix of jobs changes (e.g., if low-wage workers lose jobs, the average rises even without anyone getting a raise).

\textbf{BEA (Bureau of Economic Analysis)} --- A federal agency within the Commerce Department that produces official statistics on the U.S. economy, including GDP, personal income, and international trade data.

\textbf{business cycle} --- The recurring pattern of expansion and contraction in economic activity over time, typically measured by changes in real GDP. A full cycle includes periods of growth followed by recession.

\textbf{capital flows} --- The movement of money across national borders for investment purposes, such as foreigners buying U.S. stocks or Americans investing in overseas factories.

\textbf{capital gains} --- The profit earned when an asset (like a stock or house) is sold for more than its purchase price. Capital gains are typically taxed at lower rates than regular income.

\textbf{composition effect} --- A statistical phenomenon where changes in average wages can be influenced by shifts in who is working, not just changes in individual pay. For example, if many low-wage workers lose their jobs during a recession, the average wage may rise even though no one received a raise---simply because the remaining workforce skews toward higher earners.

\textbf{consumer spending} --- The total amount of money households spend on goods and services. It is the largest component of GDP in the United States, typically accounting for about two-thirds of economic activity.

\textbf{contribution to GDP growth} --- The portion of overall GDP growth attributable to a specific component (like consumer spending or business investment). If GDP grew 3\% and consumer spending contributed 2 percentage points, that means consumer spending alone would have produced 2\% growth.

\textbf{core inflation} --- A measure of inflation that excludes food and energy prices, which tend to be volatile. Economists use core inflation to identify underlying price trends without the noise of short-term swings in gas or grocery prices.

\textbf{corporate profits} --- The earnings companies have left after paying their costs, including wages, materials, and interest. This measure helps gauge business health and can signal future hiring or investment.

\textbf{CPI (Consumer Price Index)} --- A measure of the average change over time in the prices paid by consumers for a basket of common goods and services, such as food, housing, and transportation. It is one of the most widely used measures of inflation.

\textbf{CPS (Current Population Survey)} --- A monthly survey of about 60,000 households conducted by the Census Bureau and Bureau of Labor Statistics. It provides the official unemployment rate and other labor market data based on asking people directly about their work status.

\textbf{current account (international)} --- A measure of a country's transactions with the rest of the world, including trade in goods and services, investment income, and transfers like remittances. When outflows exceed inflows, the country runs a current account deficit. The U.S. current account deficit is driven primarily by its trade deficit, but also reflects income paid to foreign investors.

\textbf{debt-to-income ratio} --- The percentage of a person's or household's income that goes toward paying debts. A ratio of 30\% means 30 cents of every dollar earned goes to debt payments.

\textbf{deflated} --- Adjusted to remove the effects of inflation, allowing for meaningful comparisons of economic values across different time periods. Also referred to as being expressed in ``real'' or ``constant'' dollars.

\textbf{depreciation} --- The decline in value of an asset over time due to wear, age, or obsolescence. In the context of currency, it refers to a decrease in a currency's value relative to other currencies.

\textbf{discretionary spending} --- Government spending that Congress decides on each year through the appropriations process, as opposed to mandatory spending (like Social Security) that continues automatically. It includes defense, education, and transportation.

\textbf{disposable personal income} --- The amount of money households have available for spending or saving after paying taxes. It is calculated as total personal income minus personal taxes.

\textbf{durable goods} --- Products designed to last three years or more, such as cars, appliances, and furniture. Spending on durable goods tends to fluctuate more than spending on other goods during the business cycle.

\textbf{ECI (Employment Cost Index)} --- A quarterly measure of how much employers pay for labor, including both wages and benefits. Unlike average hourly earnings, the ECI holds the mix of jobs constant, giving a cleaner picture of whether compensation is actually rising.

\textbf{economic expansion} --- A period when the economy is growing---GDP is rising, jobs are being added, and business activity is increasing. Expansions end when a recession begins.

\textbf{economic growth} --- An increase in the production of goods and services in an economy over time, typically measured as the percentage change in real GDP.

\textbf{employment rate} --- The percentage of the working-age population that is employed. It is calculated by dividing the number of employed people by the total working-age population.

\textbf{employment-to-population ratio} --- The percentage of the working-age population (usually ages 16 and older) that has a job. Unlike the unemployment rate, this measure captures people who have stopped looking for work.

\textbf{equity} --- Ownership value in an asset after subtracting what is owed. For a home worth \$400,000 with a \$250,000 mortgage, the equity is \$150,000. For stocks, equity represents ownership shares in a company.

\textbf{establishment survey} --- A monthly survey of about 670,000 business locations that provides data on jobs, hours, and earnings. It produces the widely reported nonfarm payrolls number and is considered more reliable for job counts than the household survey.

\textbf{exchange rate} --- The price of one country's currency expressed in terms of another country's currency. For example, if one U.S. dollar equals 0.85 euros, that is the dollar-to-euro exchange rate.

\textbf{Fed Funds Rate} --- The interest rate banks charge each other for overnight loans. The Federal Reserve uses this rate as its main tool for influencing the economy---lowering it to encourage borrowing and spending, raising it to cool inflation.

\textbf{Federal Reserve} --- The central bank of the United States, responsible for conducting monetary policy, supervising banks, and maintaining financial stability. Often called ``the Fed.''

\textbf{financial accounts} --- A comprehensive record of the assets, liabilities, and net worth of all sectors in the economy---households, businesses, government, and the rest of the world. These accounts track who owns what and who owes what to whom. Also called the ``flow of funds.''

\textbf{fiscal stimulus} --- Government actions to boost the economy through increased spending or tax cuts. During recessions, stimulus puts money in people's pockets to encourage spending and support businesses.

\textbf{fiscal year} --- A 12-month period used for government accounting and budgeting purposes. The U.S. federal fiscal year runs from October 1 through September 30.

\textbf{FOMC (Federal Open Market Committee)} --- The Federal Reserve committee that sets interest rate policy. It meets eight times a year and includes the Fed's Board of Governors plus regional Fed bank presidents.

\textbf{GDP (Gross Domestic Product)} --- The total market value of all goods and services produced within a country during a specific period, usually a year or quarter. It is the most commonly used measure of an economy's size.

\textbf{GDPNow} --- A real-time estimate of U.S. economic growth published by the Federal Reserve Bank of Atlanta. Unlike official GDP figures, which are released quarterly with a delay, GDPNow updates frequently as new economic data becomes available, giving an early read on how the economy is performing.

\textbf{government deficit} --- The amount by which government spending exceeds government revenue in a given period. When the government collects more than it spends, it has a surplus.

\textbf{Great Recession} --- The severe economic downturn from December 2007 to June 2009, triggered by the housing market collapse and financial crisis. It was the worst U.S. recession since the 1930s, with unemployment peaking at 10\%.

\textbf{gross domestic income (GDI)} --- The total income earned producing goods and services in a country---wages, profits, and other income combined. In theory, GDI should equal GDP since every dollar spent becomes someone's income, but measurement differences cause small gaps.

\textbf{home equity} --- The portion of a home's value that the owner actually owns---the market value minus the remaining mortgage balance. Home equity often represents the largest asset for middle-class families.

\textbf{homeownership rate} --- The percentage of households that own their home rather than rent. The U.S. rate has historically hovered around 65\%.

\textbf{household survey} --- The monthly survey of individuals and families (the Current Population Survey) that produces the official unemployment rate by asking people directly whether they are working or looking for work.

\textbf{income inequality} --- The extent to which income is distributed unevenly among a population. Higher inequality means a larger gap between what high earners and low earners receive.

\textbf{inflation} --- A general increase in prices across the economy over time, which reduces the purchasing power of money. When inflation occurs, each dollar buys fewer goods and services than before.

\textbf{inflation expectations} --- What people and businesses think inflation will be in the future. These expectations matter because they influence wage negotiations, pricing decisions, and interest rates.

\textbf{inflation rate} --- The percentage change in the price level over a specific period, usually measured annually. It indicates how quickly prices are rising or falling.

\textbf{initial claims} --- The number of people filing for unemployment insurance benefits for the first time each week. A leading indicator of labor market health---rising claims often signal upcoming job losses.

\textbf{insured unemployed} --- People who are currently receiving unemployment insurance benefits. This count is smaller than total unemployment because not everyone who loses a job qualifies for benefits, and benefits eventually expire for those who remain unemployed.

\textbf{inverted yield curve} --- When short-term interest rates are higher than long-term rates, the opposite of normal. An inverted yield curve has preceded every U.S. recession in recent decades, making it a closely watched warning sign.

\textbf{job openings} --- Positions that employers are actively trying to fill. High job openings relative to unemployed workers indicate a tight labor market where workers have bargaining power.

\textbf{JOLTS (Job Openings and Labor Turnover Survey)} --- A monthly survey tracking job openings, hiring, and separations (quits, layoffs, and other departures). It provides insight into labor market dynamics beyond simple job counts.

\textbf{labor force} --- The total number of people who are either employed or actively seeking employment. It does not include people who are not working and not looking for work.

\textbf{labor force participation rate} --- The percentage of the working-age population that is in the labor force (either employed or actively seeking work). It indicates how many people are engaged in or available for work.

\textbf{labor productivity} --- The amount of goods and services produced per hour of work. Higher productivity means workers are producing more output in the same amount of time.

\textbf{labor share of income} --- The portion of national income that goes to workers as wages and benefits, as opposed to capital owners as profits. A declining labor share means workers are capturing less of the economy's gains.

\textbf{liabilities} --- What a person, company, or government owes---debts and other financial obligations. On a balance sheet, assets minus liabilities equals net worth.

\textbf{liquidity} --- How easily an asset can be converted to cash without losing value. A savings account is highly liquid; a house is not, since selling takes time and involves costs.

\textbf{M2} --- A measure of the money supply that includes cash, checking accounts, savings accounts, and other easily accessible funds. The Federal Reserve tracks M2 to gauge how much money is available for spending in the economy.

\textbf{market capitalization} --- The total value of a company's outstanding stock, calculated by multiplying share price by number of shares. A company with 1 million shares at \$100 each has a market cap of \$100 million.

\textbf{means-tested benefits} --- Government programs that only provide assistance to people below certain income or asset thresholds, such as Medicaid, food stamps (SNAP), and housing assistance.

\textbf{median} --- The middle value in a set of numbers arranged in order. Half of all values fall above the median and half fall below. Unlike an average, the median is not skewed by extremely high or low values.

\textbf{monetary policy} --- Actions taken by the Federal Reserve to influence the availability and cost of money and credit in the economy. The Fed's primary tools include setting interest rates and buying or selling securities.

\textbf{net worth} --- The total value of assets (such as homes, savings, and investments) minus total liabilities (such as mortgages and debts). It represents overall financial wealth.

\textbf{NIPA (National Income and Product Accounts)} --- The official accounting system for the U.S. economy, maintained by the BEA. It includes GDP, national income, personal income, and related measures.

\textbf{nominal GDP} --- The total value of goods and services produced in an economy measured at current prices, without adjusting for inflation. Comparing nominal GDP across years can be misleading because price changes are included.

\textbf{nonfarm payrolls} --- The total number of paid workers in the U.S. economy, excluding farm workers, private household employees, and nonprofit workers. The monthly change in nonfarm payrolls is the most-watched indicator of job growth.

\textbf{owners' equivalent rent} --- A component of the inflation measure that estimates what homeowners would pay to rent their own homes. It's used because the costs of owning (like maintenance) are considered consumption, while the home purchase itself is investment.

\textbf{payroll taxes} --- Taxes deducted from workers' paychecks to fund Social Security and Medicare. Employers also pay a matching amount, making the combined rate about 15\% of wages up to certain limits.

\textbf{PCE (Personal Consumption Expenditures)} --- A measure of consumer spending on goods and services that is used to track price changes and calculate inflation. The Federal Reserve prefers the PCE price index over the CPI for setting monetary policy.

\textbf{percentage point} --- A unit for describing the difference between two percentages. If unemployment falls from 5\% to 4\%, it dropped by 1 percentage point (not 1\%, which would be a much smaller change to 4.95\%).

\textbf{percentile} --- A value indicating the percentage of a distribution that falls below it. For example, a household at the 75th percentile of income earns more than 75\% of all households.

\textbf{personal saving rate} --- The percentage of after-tax income that households save rather than spend. A rate of 5\% means households save 5 cents of every dollar of disposable income.

\textbf{poverty threshold} --- The income level below which a family is officially considered poor. The thresholds vary by family size and are updated yearly for inflation---in 2024, about \$31,000 for a family of four.

\textbf{PPI (Producer Price Index)} --- A measure of inflation at the wholesale level, tracking prices received by producers for their output. It's often seen as a leading indicator of consumer inflation.

\textbf{producer prices} --- The prices businesses receive for their goods and services at various stages of production. Changes in producer prices can signal future changes in consumer prices.

\textbf{productivity-pay gap} --- The growing divide between how much workers produce per hour and how much they are paid. Since the 1970s, productivity in the U.S. has risen significantly faster than typical worker compensation, meaning the economic gains from increased efficiency have not been evenly shared with workers.

\textbf{public debt} --- The total amount of money the government owes to creditors, accumulated from past deficits. Also called the national debt or government debt.

\textbf{quantitative easing} --- A Federal Reserve policy of buying large amounts of bonds to push down long-term interest rates and stimulate the economy when short-term rates are already near zero. The Fed used this tool extensively after the 2008 financial crisis and during the COVID-19 pandemic.

\textbf{quintile / income quintile} --- One of five equal groups (each representing 20\%) when a population is ranked from lowest to highest. The bottom quintile is the poorest 20\%; the top quintile is the richest 20\%. Often used to describe income distribution.

\textbf{real GDP} --- The value of goods and services produced in an economy, adjusted for inflation to reflect changes in actual output rather than price changes. It allows for meaningful comparisons across time.

\textbf{real interest rates} --- Interest rates adjusted for inflation, representing the true cost of borrowing or return on saving. If a bond pays 5\% and inflation is 3\%, the real interest rate is about 2\%.

\textbf{recession} --- A significant decline in economic activity lasting more than a few months, typically visible in falling real GDP, rising unemployment, and reduced consumer spending. In the United States, the National Bureau of Economic Research (NBER) officially determines when recessions begin and end.

\textbf{residential investment} --- Spending on housing, including new home construction, home improvements, and brokers' fees. It's a component of GDP that tends to lead the broader economy---declining before recessions and recovering before expansions.

\textbf{S\&P 500} --- A stock market index tracking 500 large U.S. companies, widely used as a benchmark for overall stock market performance. It represents about 80\% of U.S. stock market value.

\textbf{Sahm Rule} --- A recession indicator that signals an economic downturn when the three-month average unemployment rate rises at least 0.5 percentage points above its lowest point in the previous 12 months. Named after economist Claudia Sahm, this rule has reliably identified every U.S. recession since 1970.

\textbf{seasonally adjusted} --- Data that has been modified to remove predictable seasonal patterns, such as holiday shopping or summer hiring. This adjustment makes it easier to identify underlying economic trends.

\textbf{sectoral balances} --- The financial positions of the major parts of the economy---households, businesses, government, and the foreign sector---showing whether each is spending more than it earns (deficit) or earning more than it spends (surplus). Because one sector's deficit must be matched by surpluses elsewhere, these balances always sum to zero across the entire economy.

\textbf{shelter (as CPI category)} --- The housing component of the Consumer Price Index, including rent and owners' equivalent rent. Shelter is the largest single category in CPI, making up about one-third of the index.

\textbf{trade balance / trade deficit / net exports} --- The difference between what a country exports and imports. When exports exceed imports, the country has a trade surplus; when imports exceed exports, a trade deficit. In GDP calculations, this same measure is called ``net exports.'' The U.S. has run a trade deficit since the 1970s, meaning Americans buy more from abroad than they sell.

\textbf{transfer payments} --- Government payments to individuals that are not in exchange for goods or services, such as Social Security, unemployment benefits, and food assistance. Transfers redistribute income but don't directly produce GDP.

\textbf{Treasury yields} --- The interest rates on U.S. government bonds. Yields vary by maturity---short-term bills might yield 4\% while 30-year bonds yield 5\%---and serve as benchmarks for other interest rates throughout the economy.

\textbf{U3 unemployment rate} --- The official unemployment rate, measuring the percentage of the labor force that is jobless and actively looking for work. It does not count people who have given up searching or those working part-time who want full-time jobs.

\textbf{U6 unemployment rate} --- A broader measure of unemployment that includes the officially unemployed plus discouraged workers (who stopped looking) and part-time workers who want full-time work. U6 is typically several percentage points higher than the official rate.

\textbf{unemployment rate} --- The percentage of the labor force that is jobless and actively seeking work. It is one of the most closely watched indicators of economic health, reported monthly.

\textbf{unit labor costs} --- The labor cost to produce one unit of output, calculated as total labor compensation divided by total output. Rising unit labor costs can signal inflation pressure; falling costs suggest productivity gains.

\textbf{value added} --- The increase in value a business creates over and above the cost of inputs it purchases. A furniture maker's value added is the difference between what it pays for wood and what it sells furniture for. GDP is essentially the sum of value added across all businesses.

\textbf{yield} --- The income earned from an investment, usually expressed as an annual percentage of the investment's value. For bonds, the yield represents the return an investor receives from interest payments.

\textbf{yield curve} --- A chart showing interest rates on bonds of different maturities, typically Treasury securities. Normally, longer-term bonds pay higher rates; when short-term rates exceed long-term rates, the curve is inverted---a recession warning signal.

\end{multicols}

% Requires: \usepackage{multicol}

\section*{Glossary}
\addcontentsline{toc}{section}{Glossary}

\begin{multicols}{2}
\small
\setlength{\parindent}{0pt}
\setlength{\parskip}{0.5\baselineskip}

\textbf{aggregate demand} --- The total demand for all goods and services in an economy at a given time. It combines spending by consumers, businesses, the government, and foreign buyers.

\textbf{annualized / annualized rate} --- A way of expressing a short-term change (such as monthly or quarterly) as if it continued for a full year. For example, if GDP grew 0.5\% in one quarter, the annualized rate would be about 2\%---the growth you'd see if that pace held for four quarters.

\textbf{average hourly earnings (AHE)} --- The average amount workers earn per hour, calculated from employer payroll data. This measure helps track wage growth but can be misleading when the mix of jobs changes (e.g., if low-wage workers lose jobs, the average rises even without anyone getting a raise).

\textbf{BEA (Bureau of Economic Analysis)} --- A federal agency within the Commerce Department that produces official statistics on the U.S. economy, including GDP, personal income, and international trade data.

\textbf{business cycle} --- The recurring pattern of expansion and contraction in economic activity over time, typically measured by changes in real GDP. A full cycle includes periods of growth followed by recession.

\textbf{capital flows} --- The movement of money across national borders for investment purposes, such as foreigners buying U.S. stocks or Americans investing in overseas factories.

\textbf{capital gains} --- The profit earned when an asset (like a stock or house) is sold for more than its purchase price. Capital gains are typically taxed at lower rates than regular income.

\textbf{composition effect} --- A statistical phenomenon where changes in average wages can be influenced by shifts in who is working, not just changes in individual pay. For example, if many low-wage workers lose their jobs during a recession, the average wage may rise even though no one received a raise---simply because the remaining workforce skews toward higher earners.

\textbf{consumer spending} --- The total amount of money households spend on goods and services. It is the largest component of GDP in the United States, typically accounting for about two-thirds of economic activity.

\textbf{contribution to GDP growth} --- The portion of overall GDP growth attributable to a specific component (like consumer spending or business investment). If GDP grew 3\% and consumer spending contributed 2 percentage points, that means consumer spending alone would have produced 2\% growth.

\textbf{core inflation} --- A measure of inflation that excludes food and energy prices, which tend to be volatile. Economists use core inflation to identify underlying price trends without the noise of short-term swings in gas or grocery prices.

\textbf{corporate profits} --- The earnings companies have left after paying their costs, including wages, materials, and interest. This measure helps gauge business health and can signal future hiring or investment.

\textbf{CPI (Consumer Price Index)} --- A measure of the average change over time in the prices paid by consumers for a basket of common goods and services, such as food, housing, and transportation. It is one of the most widely used measures of inflation.

\textbf{CPS (Current Population Survey)} --- A monthly survey of about 60,000 households conducted by the Census Bureau and Bureau of Labor Statistics. It provides the official unemployment rate and other labor market data based on asking people directly about their work status.

\textbf{current account (international)} --- A measure of a country's transactions with the rest of the world, including trade in goods and services, investment income, and transfers like remittances. When outflows exceed inflows, the country runs a current account deficit. The U.S. current account deficit is driven primarily by its trade deficit, but also reflects income paid to foreign investors.

\textbf{debt-to-income ratio} --- The percentage of a person's or household's income that goes toward paying debts. A ratio of 30\% means 30 cents of every dollar earned goes to debt payments.

\textbf{deflated} --- Adjusted to remove the effects of inflation, allowing for meaningful comparisons of economic values across different time periods. Also referred to as being expressed in ``real'' or ``constant'' dollars.

\textbf{depreciation} --- The decline in value of an asset over time due to wear, age, or obsolescence. In the context of currency, it refers to a decrease in a currency's value relative to other currencies.

\textbf{discretionary spending} --- Government spending that Congress decides on each year through the appropriations process, as opposed to mandatory spending (like Social Security) that continues automatically. It includes defense, education, and transportation.

\textbf{disposable personal income} --- The amount of money households have available for spending or saving after paying taxes. It is calculated as total personal income minus personal taxes.

\textbf{durable goods} --- Products designed to last three years or more, such as cars, appliances, and furniture. Spending on durable goods tends to fluctuate more than spending on other goods during the business cycle.

\textbf{ECI (Employment Cost Index)} --- A quarterly measure of how much employers pay for labor, including both wages and benefits. Unlike average hourly earnings, the ECI holds the mix of jobs constant, giving a cleaner picture of whether compensation is actually rising.

\textbf{economic expansion} --- A period when the economy is growing---GDP is rising, jobs are being added, and business activity is increasing. Expansions end when a recession begins.

\textbf{economic growth} --- An increase in the production of goods and services in an economy over time, typically measured as the percentage change in real GDP.

\textbf{employment rate} --- The percentage of the working-age population that is employed. It is calculated by dividing the number of employed people by the total working-age population.

\textbf{employment-to-population ratio} --- The percentage of the working-age population (usually ages 16 and older) that has a job. Unlike the unemployment rate, this measure captures people who have stopped looking for work.

\textbf{equity} --- Ownership value in an asset after subtracting what is owed. For a home worth \$400,000 with a \$250,000 mortgage, the equity is \$150,000. For stocks, equity represents ownership shares in a company.

\textbf{establishment survey} --- A monthly survey of about 670,000 business locations that provides data on jobs, hours, and earnings. It produces the widely reported nonfarm payrolls number and is considered more reliable for job counts than the household survey.

\textbf{exchange rate} --- The price of one country's currency expressed in terms of another country's currency. For example, if one U.S. dollar equals 0.85 euros, that is the dollar-to-euro exchange rate.

\textbf{Fed Funds Rate} --- The interest rate banks charge each other for overnight loans. The Federal Reserve uses this rate as its main tool for influencing the economy---lowering it to encourage borrowing and spending, raising it to cool inflation.

\textbf{Federal Reserve} --- The central bank of the United States, responsible for conducting monetary policy, supervising banks, and maintaining financial stability. Often called ``the Fed.''

\textbf{financial accounts} --- A comprehensive record of the assets, liabilities, and net worth of all sectors in the economy---households, businesses, government, and the rest of the world. These accounts track who owns what and who owes what to whom. Also called the ``flow of funds.''

\textbf{fiscal stimulus} --- Government actions to boost the economy through increased spending or tax cuts. During recessions, stimulus puts money in people's pockets to encourage spending and support businesses.

\textbf{fiscal year} --- A 12-month period used for government accounting and budgeting purposes. The U.S. federal fiscal year runs from October 1 through September 30.

\textbf{FOMC (Federal Open Market Committee)} --- The Federal Reserve committee that sets interest rate policy. It meets eight times a year and includes the Fed's Board of Governors plus regional Fed bank presidents.

\textbf{GDP (Gross Domestic Product)} --- The total market value of all goods and services produced within a country during a specific period, usually a year or quarter. It is the most commonly used measure of an economy's size.

\textbf{GDPNow} --- A real-time estimate of U.S. economic growth published by the Federal Reserve Bank of Atlanta. Unlike official GDP figures, which are released quarterly with a delay, GDPNow updates frequently as new economic data becomes available, giving an early read on how the economy is performing.

\textbf{government deficit} --- The amount by which government spending exceeds government revenue in a given period. When the government collects more than it spends, it has a surplus.

\textbf{Great Recession} --- The severe economic downturn from December 2007 to June 2009, triggered by the housing market collapse and financial crisis. It was the worst U.S. recession since the 1930s, with unemployment peaking at 10\%.

\textbf{gross domestic income (GDI)} --- The total income earned producing goods and services in a country---wages, profits, and other income combined. In theory, GDI should equal GDP since every dollar spent becomes someone's income, but measurement differences cause small gaps.

\textbf{home equity} --- The portion of a home's value that the owner actually owns---the market value minus the remaining mortgage balance. Home equity often represents the largest asset for middle-class families.

\textbf{homeownership rate} --- The percentage of households that own their home rather than rent. The U.S. rate has historically hovered around 65\%.

\textbf{household survey} --- The monthly survey of individuals and families (the Current Population Survey) that produces the official unemployment rate by asking people directly whether they are working or looking for work.

\textbf{income inequality} --- The extent to which income is distributed unevenly among a population. Higher inequality means a larger gap between what high earners and low earners receive.

\textbf{inflation} --- A general increase in prices across the economy over time, which reduces the purchasing power of money. When inflation occurs, each dollar buys fewer goods and services than before.

\textbf{inflation expectations} --- What people and businesses think inflation will be in the future. These expectations matter because they influence wage negotiations, pricing decisions, and interest rates.

\textbf{inflation rate} --- The percentage change in the price level over a specific period, usually measured annually. It indicates how quickly prices are rising or falling.

\textbf{initial claims} --- The number of people filing for unemployment insurance benefits for the first time each week. A leading indicator of labor market health---rising claims often signal upcoming job losses.

\textbf{insured unemployed} --- People who are currently receiving unemployment insurance benefits. This count is smaller than total unemployment because not everyone who loses a job qualifies for benefits, and benefits eventually expire for those who remain unemployed.

\textbf{inverted yield curve} --- When short-term interest rates are higher than long-term rates, the opposite of normal. An inverted yield curve has preceded every U.S. recession in recent decades, making it a closely watched warning sign.

\textbf{job openings} --- Positions that employers are actively trying to fill. High job openings relative to unemployed workers indicate a tight labor market where workers have bargaining power.

\textbf{JOLTS (Job Openings and Labor Turnover Survey)} --- A monthly survey tracking job openings, hiring, and separations (quits, layoffs, and other departures). It provides insight into labor market dynamics beyond simple job counts.

\textbf{labor force} --- The total number of people who are either employed or actively seeking employment. It does not include people who are not working and not looking for work.

\textbf{labor force participation rate} --- The percentage of the working-age population that is in the labor force (either employed or actively seeking work). It indicates how many people are engaged in or available for work.

\textbf{labor productivity} --- The amount of goods and services produced per hour of work. Higher productivity means workers are producing more output in the same amount of time.

\textbf{labor share of income} --- The portion of national income that goes to workers as wages and benefits, as opposed to capital owners as profits. A declining labor share means workers are capturing less of the economy's gains.

\textbf{liabilities} --- What a person, company, or government owes---debts and other financial obligations. On a balance sheet, assets minus liabilities equals net worth.

\textbf{liquidity} --- How easily an asset can be converted to cash without losing value. A savings account is highly liquid; a house is not, since selling takes time and involves costs.

\textbf{M2} --- A measure of the money supply that includes cash, checking accounts, savings accounts, and other easily accessible funds. The Federal Reserve tracks M2 to gauge how much money is available for spending in the economy.

\textbf{market capitalization} --- The total value of a company's outstanding stock, calculated by multiplying share price by number of shares. A company with 1 million shares at \$100 each has a market cap of \$100 million.

\textbf{means-tested benefits} --- Government programs that only provide assistance to people below certain income or asset thresholds, such as Medicaid, food stamps (SNAP), and housing assistance.

\textbf{median} --- The middle value in a set of numbers arranged in order. Half of all values fall above the median and half fall below. Unlike an average, the median is not skewed by extremely high or low values.

\textbf{monetary policy} --- Actions taken by the Federal Reserve to influence the availability and cost of money and credit in the economy. The Fed's primary tools include setting interest rates and buying or selling securities.

\textbf{net worth} --- The total value of assets (such as homes, savings, and investments) minus total liabilities (such as mortgages and debts). It represents overall financial wealth.

\textbf{NIPA (National Income and Product Accounts)} --- The official accounting system for the U.S. economy, maintained by the BEA. It includes GDP, national income, personal income, and related measures.

\textbf{nominal GDP} --- The total value of goods and services produced in an economy measured at current prices, without adjusting for inflation. Comparing nominal GDP across years can be misleading because price changes are included.

\textbf{nonfarm payrolls} --- The total number of paid workers in the U.S. economy, excluding farm workers, private household employees, and nonprofit workers. The monthly change in nonfarm payrolls is the most-watched indicator of job growth.

\textbf{owners' equivalent rent} --- A component of the inflation measure that estimates what homeowners would pay to rent their own homes. It's used because the costs of owning (like maintenance) are considered consumption, while the home purchase itself is investment.

\textbf{payroll taxes} --- Taxes deducted from workers' paychecks to fund Social Security and Medicare. Employers also pay a matching amount, making the combined rate about 15\% of wages up to certain limits.

\textbf{PCE (Personal Consumption Expenditures)} --- A measure of consumer spending on goods and services that is used to track price changes and calculate inflation. The Federal Reserve prefers the PCE price index over the CPI for setting monetary policy.

\textbf{percentage point} --- A unit for describing the difference between two percentages. If unemployment falls from 5\% to 4\%, it dropped by 1 percentage point (not 1\%, which would be a much smaller change to 4.95\%).

\textbf{percentile} --- A value indicating the percentage of a distribution that falls below it. For example, a household at the 75th percentile of income earns more than 75\% of all households.

\textbf{personal saving rate} --- The percentage of after-tax income that households save rather than spend. A rate of 5\% means households save 5 cents of every dollar of disposable income.

\textbf{poverty threshold} --- The income level below which a family is officially considered poor. The thresholds vary by family size and are updated yearly for inflation---in 2024, about \$31,000 for a family of four.

\textbf{PPI (Producer Price Index)} --- A measure of inflation at the wholesale level, tracking prices received by producers for their output. It's often seen as a leading indicator of consumer inflation.

\textbf{producer prices} --- The prices businesses receive for their goods and services at various stages of production. Changes in producer prices can signal future changes in consumer prices.

\textbf{productivity-pay gap} --- The growing divide between how much workers produce per hour and how much they are paid. Since the 1970s, productivity in the U.S. has risen significantly faster than typical worker compensation, meaning the economic gains from increased efficiency have not been evenly shared with workers.

\textbf{public debt} --- The total amount of money the government owes to creditors, accumulated from past deficits. Also called the national debt or government debt.

\textbf{quantitative easing} --- A Federal Reserve policy of buying large amounts of bonds to push down long-term interest rates and stimulate the economy when short-term rates are already near zero. The Fed used this tool extensively after the 2008 financial crisis and during the COVID-19 pandemic.

\textbf{quintile / income quintile} --- One of five equal groups (each representing 20\%) when a population is ranked from lowest to highest. The bottom quintile is the poorest 20\%; the top quintile is the richest 20\%. Often used to describe income distribution.

\textbf{real GDP} --- The value of goods and services produced in an economy, adjusted for inflation to reflect changes in actual output rather than price changes. It allows for meaningful comparisons across time.

\textbf{real interest rates} --- Interest rates adjusted for inflation, representing the true cost of borrowing or return on saving. If a bond pays 5\% and inflation is 3\%, the real interest rate is about 2\%.

\textbf{recession} --- A significant decline in economic activity lasting more than a few months, typically visible in falling real GDP, rising unemployment, and reduced consumer spending. In the United States, the National Bureau of Economic Research (NBER) officially determines when recessions begin and end.

\textbf{residential investment} --- Spending on housing, including new home construction, home improvements, and brokers' fees. It's a component of GDP that tends to lead the broader economy---declining before recessions and recovering before expansions.

\textbf{S\&P 500} --- A stock market index tracking 500 large U.S. companies, widely used as a benchmark for overall stock market performance. It represents about 80\% of U.S. stock market value.

\textbf{Sahm Rule} --- A recession indicator that signals an economic downturn when the three-month average unemployment rate rises at least 0.5 percentage points above its lowest point in the previous 12 months. Named after economist Claudia Sahm, this rule has reliably identified every U.S. recession since 1970.

\textbf{seasonally adjusted} --- Data that has been modified to remove predictable seasonal patterns, such as holiday shopping or summer hiring. This adjustment makes it easier to identify underlying economic trends.

\textbf{sectoral balances} --- The financial positions of the major parts of the economy---households, businesses, government, and the foreign sector---showing whether each is spending more than it earns (deficit) or earning more than it spends (surplus). Because one sector's deficit must be matched by surpluses elsewhere, these balances always sum to zero across the entire economy.

\textbf{shelter (as CPI category)} --- The housing component of the Consumer Price Index, including rent and owners' equivalent rent. Shelter is the largest single category in CPI, making up about one-third of the index.

\textbf{trade balance / trade deficit / net exports} --- The difference between what a country exports and imports. When exports exceed imports, the country has a trade surplus; when imports exceed exports, a trade deficit. In GDP calculations, this same measure is called ``net exports.'' The U.S. has run a trade deficit since the 1970s, meaning Americans buy more from abroad than they sell.

\textbf{transfer payments} --- Government payments to individuals that are not in exchange for goods or services, such as Social Security, unemployment benefits, and food assistance. Transfers redistribute income but don't directly produce GDP.

\textbf{Treasury yields} --- The interest rates on U.S. government bonds. Yields vary by maturity---short-term bills might yield 4\% while 30-year bonds yield 5\%---and serve as benchmarks for other interest rates throughout the economy.

\textbf{U3 unemployment rate} --- The official unemployment rate, measuring the percentage of the labor force that is jobless and actively looking for work. It does not count people who have given up searching or those working part-time who want full-time jobs.

\textbf{U6 unemployment rate} --- A broader measure of unemployment that includes the officially unemployed plus discouraged workers (who stopped looking) and part-time workers who want full-time work. U6 is typically several percentage points higher than the official rate.

\textbf{unemployment rate} --- The percentage of the labor force that is jobless and actively seeking work. It is one of the most closely watched indicators of economic health, reported monthly.

\textbf{unit labor costs} --- The labor cost to produce one unit of output, calculated as total labor compensation divided by total output. Rising unit labor costs can signal inflation pressure; falling costs suggest productivity gains.

\textbf{value added} --- The increase in value a business creates over and above the cost of inputs it purchases. A furniture maker's value added is the difference between what it pays for wood and what it sells furniture for. GDP is essentially the sum of value added across all businesses.

\textbf{yield} --- The income earned from an investment, usually expressed as an annual percentage of the investment's value. For bonds, the yield represents the return an investor receives from interest payments.

\textbf{yield curve} --- A chart showing interest rates on bonds of different maturities, typically Treasury securities. Normally, longer-term bonds pay higher rates; when short-term rates exceed long-term rates, the curve is inverted---a recession warning signal.

\end{multicols}

% Requires: \usepackage{multicol}

\section*{Glossary}
\addcontentsline{toc}{section}{Glossary}

\begin{multicols}{2}
\small
\setlength{\parindent}{0pt}
\setlength{\parskip}{0.5\baselineskip}

\textbf{aggregate demand} --- The total demand for all goods and services in an economy at a given time. It combines spending by consumers, businesses, the government, and foreign buyers.

\textbf{annualized / annualized rate} --- A way of expressing a short-term change (such as monthly or quarterly) as if it continued for a full year. For example, if GDP grew 0.5\% in one quarter, the annualized rate would be about 2\%---the growth you'd see if that pace held for four quarters.

\textbf{average hourly earnings (AHE)} --- The average amount workers earn per hour, calculated from employer payroll data. This measure helps track wage growth but can be misleading when the mix of jobs changes (e.g., if low-wage workers lose jobs, the average rises even without anyone getting a raise).

\textbf{BEA (Bureau of Economic Analysis)} --- A federal agency within the Commerce Department that produces official statistics on the U.S. economy, including GDP, personal income, and international trade data.

\textbf{business cycle} --- The recurring pattern of expansion and contraction in economic activity over time, typically measured by changes in real GDP. A full cycle includes periods of growth followed by recession.

\textbf{capital flows} --- The movement of money across national borders for investment purposes, such as foreigners buying U.S. stocks or Americans investing in overseas factories.

\textbf{capital gains} --- The profit earned when an asset (like a stock or house) is sold for more than its purchase price. Capital gains are typically taxed at lower rates than regular income.

\textbf{composition effect} --- A statistical phenomenon where changes in average wages can be influenced by shifts in who is working, not just changes in individual pay. For example, if many low-wage workers lose their jobs during a recession, the average wage may rise even though no one received a raise---simply because the remaining workforce skews toward higher earners.

\textbf{consumer spending} --- The total amount of money households spend on goods and services. It is the largest component of GDP in the United States, typically accounting for about two-thirds of economic activity.

\textbf{contribution to GDP growth} --- The portion of overall GDP growth attributable to a specific component (like consumer spending or business investment). If GDP grew 3\% and consumer spending contributed 2 percentage points, that means consumer spending alone would have produced 2\% growth.

\textbf{core inflation} --- A measure of inflation that excludes food and energy prices, which tend to be volatile. Economists use core inflation to identify underlying price trends without the noise of short-term swings in gas or grocery prices.

\textbf{corporate profits} --- The earnings companies have left after paying their costs, including wages, materials, and interest. This measure helps gauge business health and can signal future hiring or investment.

\textbf{CPI (Consumer Price Index)} --- A measure of the average change over time in the prices paid by consumers for a basket of common goods and services, such as food, housing, and transportation. It is one of the most widely used measures of inflation.

\textbf{CPS (Current Population Survey)} --- A monthly survey of about 60,000 households conducted by the Census Bureau and Bureau of Labor Statistics. It provides the official unemployment rate and other labor market data based on asking people directly about their work status.

\textbf{current account (international)} --- A measure of a country's transactions with the rest of the world, including trade in goods and services, investment income, and transfers like remittances. When outflows exceed inflows, the country runs a current account deficit. The U.S. current account deficit is driven primarily by its trade deficit, but also reflects income paid to foreign investors.

\textbf{debt-to-income ratio} --- The percentage of a person's or household's income that goes toward paying debts. A ratio of 30\% means 30 cents of every dollar earned goes to debt payments.

\textbf{deflated} --- Adjusted to remove the effects of inflation, allowing for meaningful comparisons of economic values across different time periods. Also referred to as being expressed in ``real'' or ``constant'' dollars.

\textbf{depreciation} --- The decline in value of an asset over time due to wear, age, or obsolescence. In the context of currency, it refers to a decrease in a currency's value relative to other currencies.

\textbf{discretionary spending} --- Government spending that Congress decides on each year through the appropriations process, as opposed to mandatory spending (like Social Security) that continues automatically. It includes defense, education, and transportation.

\textbf{disposable personal income} --- The amount of money households have available for spending or saving after paying taxes. It is calculated as total personal income minus personal taxes.

\textbf{durable goods} --- Products designed to last three years or more, such as cars, appliances, and furniture. Spending on durable goods tends to fluctuate more than spending on other goods during the business cycle.

\textbf{ECI (Employment Cost Index)} --- A quarterly measure of how much employers pay for labor, including both wages and benefits. Unlike average hourly earnings, the ECI holds the mix of jobs constant, giving a cleaner picture of whether compensation is actually rising.

\textbf{economic expansion} --- A period when the economy is growing---GDP is rising, jobs are being added, and business activity is increasing. Expansions end when a recession begins.

\textbf{economic growth} --- An increase in the production of goods and services in an economy over time, typically measured as the percentage change in real GDP.

\textbf{employment rate} --- The percentage of the working-age population that is employed. It is calculated by dividing the number of employed people by the total working-age population.

\textbf{employment-to-population ratio} --- The percentage of the working-age population (usually ages 16 and older) that has a job. Unlike the unemployment rate, this measure captures people who have stopped looking for work.

\textbf{equity} --- Ownership value in an asset after subtracting what is owed. For a home worth \$400,000 with a \$250,000 mortgage, the equity is \$150,000. For stocks, equity represents ownership shares in a company.

\textbf{establishment survey} --- A monthly survey of about 670,000 business locations that provides data on jobs, hours, and earnings. It produces the widely reported nonfarm payrolls number and is considered more reliable for job counts than the household survey.

\textbf{exchange rate} --- The price of one country's currency expressed in terms of another country's currency. For example, if one U.S. dollar equals 0.85 euros, that is the dollar-to-euro exchange rate.

\textbf{Fed Funds Rate} --- The interest rate banks charge each other for overnight loans. The Federal Reserve uses this rate as its main tool for influencing the economy---lowering it to encourage borrowing and spending, raising it to cool inflation.

\textbf{Federal Reserve} --- The central bank of the United States, responsible for conducting monetary policy, supervising banks, and maintaining financial stability. Often called ``the Fed.''

\textbf{financial accounts} --- A comprehensive record of the assets, liabilities, and net worth of all sectors in the economy---households, businesses, government, and the rest of the world. These accounts track who owns what and who owes what to whom. Also called the ``flow of funds.''

\textbf{fiscal stimulus} --- Government actions to boost the economy through increased spending or tax cuts. During recessions, stimulus puts money in people's pockets to encourage spending and support businesses.

\textbf{fiscal year} --- A 12-month period used for government accounting and budgeting purposes. The U.S. federal fiscal year runs from October 1 through September 30.

\textbf{FOMC (Federal Open Market Committee)} --- The Federal Reserve committee that sets interest rate policy. It meets eight times a year and includes the Fed's Board of Governors plus regional Fed bank presidents.

\textbf{GDP (Gross Domestic Product)} --- The total market value of all goods and services produced within a country during a specific period, usually a year or quarter. It is the most commonly used measure of an economy's size.

\textbf{GDPNow} --- A real-time estimate of U.S. economic growth published by the Federal Reserve Bank of Atlanta. Unlike official GDP figures, which are released quarterly with a delay, GDPNow updates frequently as new economic data becomes available, giving an early read on how the economy is performing.

\textbf{government deficit} --- The amount by which government spending exceeds government revenue in a given period. When the government collects more than it spends, it has a surplus.

\textbf{Great Recession} --- The severe economic downturn from December 2007 to June 2009, triggered by the housing market collapse and financial crisis. It was the worst U.S. recession since the 1930s, with unemployment peaking at 10\%.

\textbf{gross domestic income (GDI)} --- The total income earned producing goods and services in a country---wages, profits, and other income combined. In theory, GDI should equal GDP since every dollar spent becomes someone's income, but measurement differences cause small gaps.

\textbf{home equity} --- The portion of a home's value that the owner actually owns---the market value minus the remaining mortgage balance. Home equity often represents the largest asset for middle-class families.

\textbf{homeownership rate} --- The percentage of households that own their home rather than rent. The U.S. rate has historically hovered around 65\%.

\textbf{household survey} --- The monthly survey of individuals and families (the Current Population Survey) that produces the official unemployment rate by asking people directly whether they are working or looking for work.

\textbf{income inequality} --- The extent to which income is distributed unevenly among a population. Higher inequality means a larger gap between what high earners and low earners receive.

\textbf{inflation} --- A general increase in prices across the economy over time, which reduces the purchasing power of money. When inflation occurs, each dollar buys fewer goods and services than before.

\textbf{inflation expectations} --- What people and businesses think inflation will be in the future. These expectations matter because they influence wage negotiations, pricing decisions, and interest rates.

\textbf{inflation rate} --- The percentage change in the price level over a specific period, usually measured annually. It indicates how quickly prices are rising or falling.

\textbf{initial claims} --- The number of people filing for unemployment insurance benefits for the first time each week. A leading indicator of labor market health---rising claims often signal upcoming job losses.

\textbf{insured unemployed} --- People who are currently receiving unemployment insurance benefits. This count is smaller than total unemployment because not everyone who loses a job qualifies for benefits, and benefits eventually expire for those who remain unemployed.

\textbf{inverted yield curve} --- When short-term interest rates are higher than long-term rates, the opposite of normal. An inverted yield curve has preceded every U.S. recession in recent decades, making it a closely watched warning sign.

\textbf{job openings} --- Positions that employers are actively trying to fill. High job openings relative to unemployed workers indicate a tight labor market where workers have bargaining power.

\textbf{JOLTS (Job Openings and Labor Turnover Survey)} --- A monthly survey tracking job openings, hiring, and separations (quits, layoffs, and other departures). It provides insight into labor market dynamics beyond simple job counts.

\textbf{labor force} --- The total number of people who are either employed or actively seeking employment. It does not include people who are not working and not looking for work.

\textbf{labor force participation rate} --- The percentage of the working-age population that is in the labor force (either employed or actively seeking work). It indicates how many people are engaged in or available for work.

\textbf{labor productivity} --- The amount of goods and services produced per hour of work. Higher productivity means workers are producing more output in the same amount of time.

\textbf{labor share of income} --- The portion of national income that goes to workers as wages and benefits, as opposed to capital owners as profits. A declining labor share means workers are capturing less of the economy's gains.

\textbf{liabilities} --- What a person, company, or government owes---debts and other financial obligations. On a balance sheet, assets minus liabilities equals net worth.

\textbf{liquidity} --- How easily an asset can be converted to cash without losing value. A savings account is highly liquid; a house is not, since selling takes time and involves costs.

\textbf{M2} --- A measure of the money supply that includes cash, checking accounts, savings accounts, and other easily accessible funds. The Federal Reserve tracks M2 to gauge how much money is available for spending in the economy.

\textbf{market capitalization} --- The total value of a company's outstanding stock, calculated by multiplying share price by number of shares. A company with 1 million shares at \$100 each has a market cap of \$100 million.

\textbf{means-tested benefits} --- Government programs that only provide assistance to people below certain income or asset thresholds, such as Medicaid, food stamps (SNAP), and housing assistance.

\textbf{median} --- The middle value in a set of numbers arranged in order. Half of all values fall above the median and half fall below. Unlike an average, the median is not skewed by extremely high or low values.

\textbf{monetary policy} --- Actions taken by the Federal Reserve to influence the availability and cost of money and credit in the economy. The Fed's primary tools include setting interest rates and buying or selling securities.

\textbf{net worth} --- The total value of assets (such as homes, savings, and investments) minus total liabilities (such as mortgages and debts). It represents overall financial wealth.

\textbf{NIPA (National Income and Product Accounts)} --- The official accounting system for the U.S. economy, maintained by the BEA. It includes GDP, national income, personal income, and related measures.

\textbf{nominal GDP} --- The total value of goods and services produced in an economy measured at current prices, without adjusting for inflation. Comparing nominal GDP across years can be misleading because price changes are included.

\textbf{nonfarm payrolls} --- The total number of paid workers in the U.S. economy, excluding farm workers, private household employees, and nonprofit workers. The monthly change in nonfarm payrolls is the most-watched indicator of job growth.

\textbf{owners' equivalent rent} --- A component of the inflation measure that estimates what homeowners would pay to rent their own homes. It's used because the costs of owning (like maintenance) are considered consumption, while the home purchase itself is investment.

\textbf{payroll taxes} --- Taxes deducted from workers' paychecks to fund Social Security and Medicare. Employers also pay a matching amount, making the combined rate about 15\% of wages up to certain limits.

\textbf{PCE (Personal Consumption Expenditures)} --- A measure of consumer spending on goods and services that is used to track price changes and calculate inflation. The Federal Reserve prefers the PCE price index over the CPI for setting monetary policy.

\textbf{percentage point} --- A unit for describing the difference between two percentages. If unemployment falls from 5\% to 4\%, it dropped by 1 percentage point (not 1\%, which would be a much smaller change to 4.95\%).

\textbf{percentile} --- A value indicating the percentage of a distribution that falls below it. For example, a household at the 75th percentile of income earns more than 75\% of all households.

\textbf{personal saving rate} --- The percentage of after-tax income that households save rather than spend. A rate of 5\% means households save 5 cents of every dollar of disposable income.

\textbf{poverty threshold} --- The income level below which a family is officially considered poor. The thresholds vary by family size and are updated yearly for inflation---in 2024, about \$31,000 for a family of four.

\textbf{PPI (Producer Price Index)} --- A measure of inflation at the wholesale level, tracking prices received by producers for their output. It's often seen as a leading indicator of consumer inflation.

\textbf{producer prices} --- The prices businesses receive for their goods and services at various stages of production. Changes in producer prices can signal future changes in consumer prices.

\textbf{productivity-pay gap} --- The growing divide between how much workers produce per hour and how much they are paid. Since the 1970s, productivity in the U.S. has risen significantly faster than typical worker compensation, meaning the economic gains from increased efficiency have not been evenly shared with workers.

\textbf{public debt} --- The total amount of money the government owes to creditors, accumulated from past deficits. Also called the national debt or government debt.

\textbf{quantitative easing} --- A Federal Reserve policy of buying large amounts of bonds to push down long-term interest rates and stimulate the economy when short-term rates are already near zero. The Fed used this tool extensively after the 2008 financial crisis and during the COVID-19 pandemic.

\textbf{quintile / income quintile} --- One of five equal groups (each representing 20\%) when a population is ranked from lowest to highest. The bottom quintile is the poorest 20\%; the top quintile is the richest 20\%. Often used to describe income distribution.

\textbf{real GDP} --- The value of goods and services produced in an economy, adjusted for inflation to reflect changes in actual output rather than price changes. It allows for meaningful comparisons across time.

\textbf{real interest rates} --- Interest rates adjusted for inflation, representing the true cost of borrowing or return on saving. If a bond pays 5\% and inflation is 3\%, the real interest rate is about 2\%.

\textbf{recession} --- A significant decline in economic activity lasting more than a few months, typically visible in falling real GDP, rising unemployment, and reduced consumer spending. In the United States, the National Bureau of Economic Research (NBER) officially determines when recessions begin and end.

\textbf{residential investment} --- Spending on housing, including new home construction, home improvements, and brokers' fees. It's a component of GDP that tends to lead the broader economy---declining before recessions and recovering before expansions.

\textbf{S\&P 500} --- A stock market index tracking 500 large U.S. companies, widely used as a benchmark for overall stock market performance. It represents about 80\% of U.S. stock market value.

\textbf{Sahm Rule} --- A recession indicator that signals an economic downturn when the three-month average unemployment rate rises at least 0.5 percentage points above its lowest point in the previous 12 months. Named after economist Claudia Sahm, this rule has reliably identified every U.S. recession since 1970.

\textbf{seasonally adjusted} --- Data that has been modified to remove predictable seasonal patterns, such as holiday shopping or summer hiring. This adjustment makes it easier to identify underlying economic trends.

\textbf{sectoral balances} --- The financial positions of the major parts of the economy---households, businesses, government, and the foreign sector---showing whether each is spending more than it earns (deficit) or earning more than it spends (surplus). Because one sector's deficit must be matched by surpluses elsewhere, these balances always sum to zero across the entire economy.

\textbf{shelter (as CPI category)} --- The housing component of the Consumer Price Index, including rent and owners' equivalent rent. Shelter is the largest single category in CPI, making up about one-third of the index.

\textbf{trade balance / trade deficit / net exports} --- The difference between what a country exports and imports. When exports exceed imports, the country has a trade surplus; when imports exceed exports, a trade deficit. In GDP calculations, this same measure is called ``net exports.'' The U.S. has run a trade deficit since the 1970s, meaning Americans buy more from abroad than they sell.

\textbf{transfer payments} --- Government payments to individuals that are not in exchange for goods or services, such as Social Security, unemployment benefits, and food assistance. Transfers redistribute income but don't directly produce GDP.

\textbf{Treasury yields} --- The interest rates on U.S. government bonds. Yields vary by maturity---short-term bills might yield 4\% while 30-year bonds yield 5\%---and serve as benchmarks for other interest rates throughout the economy.

\textbf{U3 unemployment rate} --- The official unemployment rate, measuring the percentage of the labor force that is jobless and actively looking for work. It does not count people who have given up searching or those working part-time who want full-time jobs.

\textbf{U6 unemployment rate} --- A broader measure of unemployment that includes the officially unemployed plus discouraged workers (who stopped looking) and part-time workers who want full-time work. U6 is typically several percentage points higher than the official rate.

\textbf{unemployment rate} --- The percentage of the labor force that is jobless and actively seeking work. It is one of the most closely watched indicators of economic health, reported monthly.

\textbf{unit labor costs} --- The labor cost to produce one unit of output, calculated as total labor compensation divided by total output. Rising unit labor costs can signal inflation pressure; falling costs suggest productivity gains.

\textbf{value added} --- The increase in value a business creates over and above the cost of inputs it purchases. A furniture maker's value added is the difference between what it pays for wood and what it sells furniture for. GDP is essentially the sum of value added across all businesses.

\textbf{yield} --- The income earned from an investment, usually expressed as an annual percentage of the investment's value. For bonds, the yield represents the return an investor receives from interest payments.

\textbf{yield curve} --- A chart showing interest rates on bonds of different maturities, typically Treasury securities. Normally, longer-term bonds pay higher rates; when short-term rates exceed long-term rates, the curve is inverted---a recession warning signal.

\end{multicols}
